\section{Wnioski}
\label{sec:wnioski}

Zaprojektowano i uruchomiono w programie LOGO!Soft Comfort
sekwencyjny układ drabinkowy obsługujący przycisk zasilania telefonu.
Jeden przycisk steruje dwoma wyjściami, a rodzaj akcji zależy od czasu
przytrzymania oraz od bieżącego stanu urządzenia.

\begin{enumerate}
    \item Poprawny opis układu wymaga pięciu stanów. Stan menu
          otwartego trzeba rozdzielić na dwa warianty, aby po
          anulowaniu menu przywrócić poprzedni stan ekranu.
    \item Najistotniejszym problemem praktycznym okazała się
          cykliczna praca sterownika. Przy przytrzymanym przycisku
          sygnał pozostaje wysoki przez wiele cykli, a timery TON
          zerują się w tym samym cyklu, w którym przycisk jest
          puszczany. Pierwotne rozwiązanie oparte na impulsie zbocza
          opadającego nie rozpoznawało kliknięć krótkich i średnich,
          a kliknięcie długie powodowało migotanie wyjścia.
    \item Problem rozwiązano przez zapamiętanie sygnałów $I_1$,
          $T_{001}$, $T_{002}$ z poprzedniego cyklu w markerach
          $M_7$--$M_9$ i wyznaczenie zdarzeń K, Ś i D jako impulsów
          jednocyklowych. Dzięki temu każde naciśnięcie przycisku
          daje dokładnie jedno zdarzenie, niezależnie od czasu
          trzymania.
    \item Warunki Set i Reset wszystkich przerzutników RS są
          rozłączne, więc poprawność układu nie zależy od priorytetu
          wejść bloku RS. Weryfikacja równań względem tabeli
          przejść (tab.~\ref{tab:weryfikacja}) potwierdziła zgodność
          dla wszystkich stanów i zdarzeń.
    \item Warunek $M_1$ w równaniu wyjścia $Q_1$ oraz zerowanie
          $M_2$ i $M_3$ przy wyłączaniu telefonu gwarantują, że
          ekran nie świeci przy wyłączonym telefonie.
    \item Ograniczeniami rozwiązania są stałe progi czasowe
          (\SI{0,5}{\second} i \SI{2}{\second}) oraz zależność od
          kolejności sieci w programie -- zmiana kolejności, zwłaszcza
          sieci~7, naruszyłaby jednocyklowy charakter zdarzeń.
\end{enumerate}
